\section{Other task models}\label{app:othermodels}
These comparisons are exploratory.
\paragraph{Claude Haiku~4.5.} This model rarely reaches the seed's cap ($\vkMaxTok\%$ of the seed's responses stop at
it, among runs with logs), but another failure takes its place: in $\vsigKSeed\%$ of the seed's runs the agent process is killed by a signal
before it finishes (\vsigKSeedTr\% in the first-round transfer runs; \vsigHSeed\% under Haiku~5.5). The chat logs
of these runs are empty, because the seed writes its log on exit; in the SWE-bench pilot with the same model, such
endings came from cleanup commands that the agent itself ran (Appendix~\ref{app:swe}). A batch of single textbook moves under this
model, in which the seed scores $\vkgNoneF\%$, separates the causes. The non-interactive shell B removes the signal endings
entirely (\vkgBSig\% of runs, compared with \vkgNoneSig\% for the seed) and gains $\vkgBgain$ points [$\vkgBgainlo$,
$\vkgBgainhi$]; the editor, the model loop and the workflow prompt alone show no gain ($\vkgEgain$, $\vkgLgain$ and
$\vkgPgain$ points), and their runs still end on a signal in about three quarters of cases. All four moves together gain
$\vkgEBLPgain$ points [$\vkgEBLPgainlo$, $\vkgEBLPgainhi$], $\vkgBMinusEBLPAbs$ more than the shell alone. Of the two
loops, one found nothing ($\vkOnegain$ [$\vkOnegainlo$, $\vkOnegainhi$] in the final batch; \vsigKOne\% of its runs end
on a signal). The other gained $\vkTwogain$ points [$\vkTwogainlo$, $\vkTwogainhi$] at \$$\vkTwoUsd$ per thousand tasks
with four edits that rebuild parts of the engineered seed's workflow (repeated attempts with the previous failure in
context, a brute-force reference with random stress tests, timing on a maximum-size input, and a rule never to stop with
a known time-out); \vsigKTwo\% of its runs end on a signal, and others reach the budget or the time limit
(Appendix~\ref{app:cens}). The seed resampled eight times, at a lower cost than that agent, is $\vkTwoMatchedAbs$ points
[$\vkTwoMatchedAbslo$, $\vkTwoMatchedAbshi$] below it on $F_1$, so its gain is not a compute effect alone.
The engineered seed, which combines the shell fix with the workflow, is $\vkStrongMinusTwo$ points
[$\vkStrongMinusTwolo$, $\vkStrongMinusTwohi$] ahead of it at a lower cost in the final batch, where the seed scores $\vkSeedF\%$. The Haiku~4.5
placebo loop's best selection score (34\%) lies between those of the two loops (28\% and 40\%).

\paragraph{gpt-oss-120b.} This model is slow: depending on the agent, \vtoOssMin{} to \vtoOssMax\% of its runs reach
the 600\,s limit, and the solution present at that point is graded (Appendix~\ref{app:cens}). In the final batch for
this model the DGM seed solves $\voSeedF\%$ of the fresh problems; the engineered seed scores $\voStronggainAbs$ points
below it. One of three loops gained $\voOnegain$ points [$\voOnegainlo$, $\voOnegainhi$], with a workflow prompt and the
rule to write the solution file at once and never end a turn without it; its best agent costs about three times as much
per task as the seed, and the seed resampled \vossSkK{} times at a similar cost reaches $\vossSkAcc\%$ on $F_1$, and the
contrast of the loop agent with it there is $\vossLoopMinusSk$ points [$\vossLoopMinusSklo$, $\vossLoopMinusSkhi$]. A
second loop's reported agent has a fresh gain of $\voThreegain$ [$\voThreegainlo$, $\voThreegainhi$] and reaches the 600\,s limit in
\vtoOssThree\% of its runs, and the third never got a child through the
full selection set; neither did any of the placebo loop's 40 copies. In a separate batch
of single textbook moves, in which the seed scores $\vogNoneF\%$, the moves as written for Haiku~5.5 do not transfer: the editor alone
scores $\vogEgainAbs$ points below the seed, the move chosen on $F_1$ (\vogWinner) gains $\vogLgain$ [$\vogLgainlo$,
$\vogLgainhi$], and the loop's agent, measured in the same batch, is ahead of that move by $\vogWinnerFtwoAbs$
[$\vogWinnerFtwoAbslo$, $\vogWinnerFtwoAbshi$] on $F_2$. Loop 3's agent from the first round, evolved with Haiku~5.5,
gained $\ntrOssthree$ points [$\ntrOssthreelo$, $\ntrOssthreehi$] on $F_1$ in the first-round transfer runs
(Appendix~\ref{app:round1}), more than any textbook move gains in the batch above; the two numbers come from different
batches. This is consistent with the retries after empty turns that the loop found transferring where our textbook set
does not, but we did not ablate them under this model.
In a third batch we evaluated five one-shot redesigns made from five failure logs of this model. The best of five,
chosen on $F_1$, beats the loop's agent on $F_2$ by $\vorBestVsLoop$ points [$\vorBestVsLooplo$, $\vorBestVsLoophi$]; the average of the five
trails it by $\vorMeanVsLoopAbs$ [$\vorMeanVsLoopAbslo$, $\vorMeanVsLoopAbshi$], because one redesign was broken and
solved nothing. Choosing the best of five used 116 evaluated problems, more than the loops' 50-problem selection set.
